\documentclass[11pt,a4paper]{article}
\usepackage{amsmath,amssymb,amsthm}
\usepackage{hyperref}
\usepackage{doi}
\usepackage{geometry}
\geometry{margin=2.5cm}

\newtheorem{theorem}{Theorem}[section]
\newtheorem{lemma}[theorem]{Lemma}
\newtheorem{corollary}[theorem]{Corollary}
\newtheorem{proposition}[theorem]{Proposition}
\newtheorem{definition}[theorem]{Definition}
\newtheorem{remark}[theorem]{Remark}

\title{Fibrewise Born--Infeld Capacity Lapse\\
and Gravitational Time Dilation\\
\large (LC-O2 of the Lorentz Capacity Programme)}
\author{J\'er\^ome Beau\\
\small Independent Researcher\\
\small \texttt{jerome.beau@cosmochrony.org}}
\date{2026}

\begin{document}

  \maketitle

  \begin{abstract}
    The note~\cite{Beau2026n} identifies the construction of a temporal residual
    map~$\mathcal{P}_\tau^{\mathrm{Q5b}}$ as the remaining step needed to close the Lorentz
    identification from O7 observables.
    The paper~\cite{Beau2026tp} carries out this construction in the homogeneous regime and
    establishes the uniqueness of the resulting map.
    The present paper extends this construction to the non-homogeneous fibrewise regime;
    the homogeneous case is recovered as Corollary~\ref{cor:homogeneous}.

    In a non-homogeneous configuration a localised spectral occupancy $\epsilon(x)$
    reduces the available BI capacity budget at each point $x$ to a local radius
    $R(x)^2 = 1-\epsilon(x)$.
    The Weil filling law and the Lorentzian splitting are unchanged.
    A fibre-by-fibre uniqueness argument forces the
    \emph{fibrewise} temporal residual map:
    \[
      \mathcal{P}_{\tau,x}^{\mathrm{Q5b}}\Sigma^{\mathrm{tot}}(x,n)
      \;=\; R(x)\,\sqrt{f(x,n)\bigl(2-f(x,n)\bigr)},
    \]
    where $f(x,n) = 1-\beta(x,n)$ and $\beta(x,n)$ is the local normalised spatial
    mobility.
    Gravitational time dilation is identified as the position-dependence of the temporal
    residual capacity:
    \[
      \frac{d\tau}{dn}\bigg|_x = R(x)\,\sqrt{1-\beta^2} = \frac{R(x)}{\gamma}.
    \]
    The local capacity radius is identified metrically as
    $R(x)^2 = -2\,g_{\tau\tau}(x)$,
    where $g_{\tau\tau}(x)$ is the inverse metric temporal component
    and the factor of $2$ comes from the flat-space reference $g_{\tau\tau}^\infty = -1/2$
    (Q11~\cite{Beau2026q11}).
    In the weak-field regime, the Newtonian profile of $g_{\tau\tau}(x)$ established
    in~\cite{Beau2026h} gives $R(x) \approx 1 + \Phi_N(x)/c^2$, recovering
    gravitational time dilation as a consequence of the local BI capacity budget.
    The open spectral problem is to derive $R(x)$ directly from the local occupancy
    $\epsilon(x)$ without importing the Einstein equations (LC-O2-O1).

    The purpose of this note is not to rederive the Einstein equation.
    The infrared Einstein response is supplied by the renormalized spectral entropy
    functional of~\cite{Beau2026h}.
    The purpose is instead to identify the local kinematic meaning of gravitational time
    dilation in the capacity-flow language:
    a non-uniform projected configuration reduces the local Born--Infeld capacity radius,
    and the temporal residual map becomes fibrewise.
  \end{abstract}

  \noindent\textbf{Keywords.}
  fibrewise Born--Infeld capacity, gravitational time dilation, capacity lapse,
  local capacity radius, temporal residual map, Lorentz factor, non-homogeneous
  effective geometry, emergent gravitation, Cosmochrony.

  \newpage
  \tableofcontents

  \newpage

  \section{Introduction}

    \subsection{From homogeneous to fibrewise}

      The note~\cite{Beau2026n} establishes, conditionally on the construction of
      $\mathcal{P}_\tau^{\mathrm{Q5b}}$, that the core identity
      \begin{equation}\label{eq:flat}
        \mathcal{P}_\tau^{\mathrm{Q5b}}\Sigma_n^{\mathrm{tot}} = \sqrt{1-B_n^2}
      \end{equation}
      recovers $\Phi(\eta_n^\tau) = 1/\gamma$.
      The paper~\cite{Beau2026tp} closes this construction unconditionally and proves
      uniqueness of the homogeneous temporal residual map.
      The present paper extends this to the fibrewise non-homogeneous regime;
      the homogeneous case ($R(x)\equiv 1$) is recovered as
      Corollary~\ref{cor:homogeneous}.
      Gravitational time dilation, in this language, asks what happens when condition~(BI)
      is modified: when the BI capacity budget is no longer a unit sphere everywhere but
      a sphere of locally varying radius $R(x) \leq 1$.

      This is the question addressed in the present paper.
      We do not ask how the local budget variation $\epsilon(x) = 1-R(x)^2$ arises dynamically;
      this is the domain of the renormalized spectral entropy functional of~\cite{Beau2026h},
      which produces the Einstein tensor as the infrared-dominant response to a
      metric-dependent spectral operator.
      We ask instead: given a position-dependent local capacity radius $R(x)$, what is
      the fibrewise extension of the temporal residual map, and what kinematic identity
      does it produce?

      The answer is structurally forced by the same uniqueness argument as in the homogeneous
      case, applied fibrewise.

    \subsection{Two choices and the canonical one}

      When the BI budget becomes local, there are two natural ways to extend the observables.

      \paragraph{Absolute capacity formulation.}
        Keep $B(x,n)$ as an absolute fraction of capacity and replace the unit sphere by
        $\bigl(F^\tau(x,n)\bigr)^2 + B(x,n)^2 = R(x)^2$.
        The temporal residual is then $F^\tau(x,n) = \sqrt{R(x)^2-(1-f(x,n))^2}$,
        defined only when $1-f(x,n) \leq R(x)$.
        This mixes the local budget into the definition of $B(x,n)$ and loses the clean
        separation between Lorentzian geometry and local scale.

      \paragraph{Normalised capacity formulation (canonical).}
        Write $B(x,n) = R(x)\,\beta(x,n)$ with $\beta(x,n) \in [0,1]$ the local normalised
        spatial mobility, and $f(x,n) = 1-\beta(x,n)$.
        The BI sphere in normalised variables is $(F^\tau)^2/R(x)^2 + \beta^2 = 1$, giving
      \begin{equation}\label{eq:Ftau-local}
        F^\tau(x,n) = R(x)\,\sqrt{1-\beta(x,n)^2}.
        \end{equation}
        The temporal residual map then factorises as
      \begin{equation}\label{eq:factorisation}
        \mathcal{P}_{\tau,x}^{\mathrm{Q5b}} \;\;\longleftrightarrow\;\;
        \underbrace{R(x)}_{\text{local budget}}
        \;\times\;
        \underbrace{\sqrt{f(x,n)(2-f(x,n))}}_{\text{normalised Lorentzian geometry}}.
        \end{equation}
        This is the canonical formulation: $R(x)$ carries all the local physics, while
        $\sqrt{f(2-f)}$ is the universal normalised Lorentzian structure
        established in~\cite{Beau2026tp}.

  \subsection{Main results and organisation}

    Section~\ref{sec:setup} fixes the fibrewise extension of the capacity bundle.
    Section~\ref{sec:spectral} defines $R(x)$ spectrally and identifies it with the metric.
    Section~\ref{sec:main} states and proves the main theorem (fibrewise map).
    Section~\ref{sec:timedil} extracts the gravitational time dilation formula.
    Section~\ref{sec:weakfield} derives the weak-field limit from~\cite{Beau2026h}.
    Section~\ref{sec:status} gives the logical status table.

  \section{Fibrewise Capacity Bundle}\label{sec:setup}

  \subsection{Notation}

    Throughout, $(M, g)$ denotes the effective Lorentzian spacetime $M = \mathbb{R}_\tau
  \times \mathrm{Heis}_3(\mathbb{R})$ with completed co-metric $g^{\mu\nu}(x)$.
    In the homogeneous regime $g^{\mu\nu} = 2\eta^{\mu\nu}$ (Q11~\cite{Beau2026q11}).
    In the non-homogeneous regime, $g^{\mu\nu}(x)$ is the effective co-metric determined
    by the infrared Einstein response of~\cite{Beau2026h}; it is Lorentzian,
    smooth in the weak-curvature regime, and reduces to $2\eta^{\mu\nu}$ at spatial infinity.

    The flat-space Casimir value is $A_\tau^{\infty} = 2$ (Q11~\cite{Beau2026q11});
    we use this as the reference value for the temporal component.

  \subsection{Fibrewise BI sphere}

    \begin{definition}[Local capacity radius]\label{def:radius}
      Let $\epsilon(x) \geq 0$ be the \emph{local spectral occupancy} at $x \in M$:
      the fraction of the BI capacity budget at $x$ already occupied by the local
      projected configuration.
      The \emph{local capacity radius} is
      \begin{equation}\label{eq:radius}
        R(x)^2 := 1 - \epsilon(x).
      \end{equation}
      In the homogeneous vacuum ($\epsilon \equiv 0$) one recovers $R(x) = 1$ and the
      unit Born--Infeld sphere $(F_n^\tau)^2 + B_n^2 = 1$.
    \end{definition}

    \begin{definition}[Local normalised mobility and fibrewise filling law]\label{def:local}
      The \emph{local normalised spatial mobility} $\beta(x,n) \in [0,1]$ is defined by
      $B(x,n) = R(x)\,\beta(x,n)$, where $B(x,n)$ is the O7 spatial coverage fraction
      at position $x$ and depth~$n$.
      The \emph{local Weil filling variable} is $f(x,n) := 1 - \beta(x,n)$.
    \end{definition}

    \begin{remark}[Fibrewise filling law]
The Weil factorization argument of~\cite{Beau2026n} (Corollary~1) gives
$\Sigma_n^{\mathrm{tot}} = 1-B_n$ globally.
In the fibrewise setting, the Weil representation is local to each BFS history,
and the same argument yields
\begin{equation}\label{eq:fillinglocal}
  \Sigma^{\mathrm{tot}}(x,n) = R(x)\,\bigl(1-\beta(x,n)\bigr) = R(x)\,f(x,n),
\end{equation}
i.e.\ the total residual is the local budget fraction not yet reached by the spatial BFS.
The linear filling law holds in the normalised variable $f = 1-\beta$:
$\beta(x,n) + f(x,n) = 1$.
    \end{remark}

    \begin{definition}[Fibrewise capacity flow vector]\label{def:fibrevector}
      The \emph{local capacity flow vector} at $(x,n)$ is
      \[
        \vec{F}(x,n) = F^\tau(x,n)\,\partial_\tau + \vec{F}^{\,\mathrm{spat}}(x,n)
        \;\in\; T_x M,
      \]
      with $|\vec{F}^{\,\mathrm{spat}}(x,n)| = B(x,n)/dn = R(x)\beta(x,n)$.
      The BI budget at $x$ places $\vec{F}(x,n)$ on the \emph{local BI sphere}:
      \begin{equation}\label{eq:BIlocal}
        \bigl(F^\tau(x,n)\bigr)^2 + \bigl(R(x)\beta(x,n)\bigr)^2 = R(x)^2.
      \end{equation}
    \end{definition}

  \section{Spectral Occupancy and Metric Identification of $R(x)$}\label{sec:spectral}

  \subsection{Local spectral occupancy from the admissibility data}

    In the Cosmochrony framework, a localised mass-energy configuration is a projected
    configuration that has depleted part of the BI capacity budget in a neighbourhood
    of $x$.
    More precisely, the admissible capacity budget at $x$ is the fraction of the spectral
    space $\mathrm{Sym}^2(V_\rho)$ available for further projective transitions:
    \begin{equation}\label{eq:epsilon}
      \epsilon(x) = \frac{\sigma_{\mathrm{loc}}(x)}{\sigma_{\max}},
    \end{equation}
    where $\sigma_{\mathrm{loc}}(x)$ is the spectral capacity already occupied by the
    configuration at~$x$ and $\sigma_{\max} = c_{\mathrm{BI}}^2$ is the maximal budget.
    In the homogeneous vacuum, $\sigma_{\mathrm{loc}} \equiv 0$ and $R(x) = 1$.

  \subsection{Identification with the effective co-metric}

    The completed effective co-metric at $x$ is $g^{\mu\nu}(x)$.
    Its temporal component satisfies $g^{\tau\tau}(x) < 0$ (Lorentzian signature).
    The flat-space reference value is $g^{\tau\tau,\infty} = -A_\tau^\infty = -2$
    (Q11~\cite{Beau2026q11}).

    \begin{definition}[Metric capacity lapse]\label{def:lapse}
      The \emph{metric capacity lapse} at $x$ is the ratio of the local temporal metric
      component to its flat-space reference value:
      \begin{equation}\label{eq:lapse}
        R(x)^2 := \frac{g_{\tau\tau}(x)}{g_{\tau\tau}^\infty}
        = -2\,g_{\tau\tau}(x),
      \end{equation}
      where $g_{\tau\tau}(x) = (g^{\tau\tau}(x))^{-1}$ is the inverse metric temporal component,
      $g_{\tau\tau}^\infty = -1/2$ is its flat-space reference value (from $g^{\mu\nu}=2\eta^{\mu\nu}$,
      Q11~\cite{Beau2026q11}), and the explicit form $-2g_{\tau\tau}(x)$ follows from
      $g_{\tau\tau}^\infty = -1/2$.
      In the homogeneous vacuum $g_{\tau\tau}^\infty = -1/2$, giving $R(\infty) = 1$.
    \end{definition}

    \begin{remark}[Postulated equivalence of the two definitions]\label{rem:equiv}
      Definitions~\ref{def:radius} and~\ref{def:lapse} are postulated to describe the same
      local capacity depletion in the Einstein regime.
      Their direct equivalence from the admissibility data alone is not proved here;
      establishing $\epsilon(x) = 1-R(x)^2 = 1+2g_{\tau\tau}(x)$ purely from spectral occupancy
      is stated as open problem LC-O2-O1.
      The metric formulation~\eqref{eq:lapse} is used for the weak-field derivation
      in Section~\ref{sec:weakfield}.
    \end{remark}

  \section{Fibrewise Temporal Residual Map}\label{sec:main}

  \subsection{The Lorentzian projector at $x$}

    The local co-metric $g^{\mu\nu}(x)$ determines a Lorentzian temporal projector
    $P_{\tau,x}$ on $T_x M$ by the same formula as in the homogeneous case
    (\cite{Beau2026tp}, Section~4.1):
    \begin{equation}\label{eq:Ptau-local}
      P_{\tau,x}\,V = V^\tau\,\partial_\tau, \qquad V \in T_x M.
    \end{equation}

  \subsection{Fibrewise temporal residual map}

    \begin{definition}[Fibrewise temporal residual map]\label{def:Ptau-local}
      The \emph{fibrewise temporal residual map} $\mathcal{P}_{\tau,x}^{\mathrm{Q5b}}$
      is the composition:
      \[
        \mathcal{P}_{\tau,x}^{\mathrm{Q5b}} \colon
        \Sigma^{\mathrm{tot}}(x,n)
        \xrightarrow{\;\text{BI lift at }x\;}
        \vec{F}(x,n)
        \xrightarrow{\;P_{\tau,x}\;}
        F^\tau(x,n)\,\partial_\tau
        \xrightarrow{\;\text{magnitude}\;}
        |F^\tau(x,n)|.
      \]
    \end{definition}

  \subsection{Main theorem}

    \begin{theorem}[Fibrewise temporal residual map]\label{thm:main}
      Assume the Q5a continuum hypotheses~\cite{Beau2026q5a},
      the spectral universality~\textup{[U]} proved in U1~\cite{Beau2026u1},
      the fibrewise BI capacity lift of Definition~\ref{def:fibrevector},
      and the local co-metric~$g^{\mu\nu}(x)$ of the effective Lorentzian geometry.
      Then the fibrewise temporal residual map acts by
      \begin{equation}\label{eq:main}
        \mathcal{P}_{\tau,x}^{\mathrm{Q5b}}\,\Sigma^{\mathrm{tot}}(x,n)
        \;=\; R(x)\,\sqrt{f(x,n)\bigl(2-f(x,n)\bigr)},
      \end{equation}
      where $f(x,n) = 1-\beta(x,n)$, $R(x)$ is the local capacity radius~\eqref{eq:lapse},
      and $\beta(x,n) \in [0,1]$ is the local normalised spatial mobility.

      The quadratic residual law holds fibrewise:
      \begin{equation}\label{eq:quadratic-local}
        \bigl(R(x)\beta(x,n)\bigr)^2
        + \bigl(\mathcal{P}_{\tau,x}^{\mathrm{Q5b}}\,\Sigma^{\mathrm{tot}}(x,n)\bigr)^2
        = R(x)^2.
      \end{equation}

      Moreover, $\mathcal{P}_{\tau,x}^{\mathrm{Q5b}}$ is the unique positive function satisfying
      simultaneously:
      \begin{itemize}
        \item[\textup{(L$_x$)}] $\mathcal{P}_{\tau,x}^{\mathrm{Q5b}}\Sigma^{\mathrm{tot}}(x,n)
        = |P_{\tau,x}\vec{F}(x,n)|$ (local Lorentzian projection);
        \item[\textup{(BI$_x$)}] $B(x,n)^2 + (\mathcal{P}_{\tau,x}^{\mathrm{Q5b}}
        \Sigma^{\mathrm{tot}}(x,n))^2 = R(x)^2$ (local BI sphere~\eqref{eq:BIlocal});
        \item[\textup{(W$_x$)}] $f(x,n) = 1-\beta(x,n)$ with
        $B(x,n) = R(x)\beta(x,n)$ (fibrewise filling law~\eqref{eq:fillinglocal}).
      \end{itemize}
    \end{theorem}

    \begin{proof}
      \textbf{Explicit formula.}
      From~(BI$_x$): $(F^\tau)^2 = R(x)^2 - B(x,n)^2 = R(x)^2 - R(x)^2\beta^2
      = R(x)^2(1-\beta^2)$.
      Hence $|F^\tau| = R(x)\sqrt{1-\beta^2}$.
      By~(W$_x$): $\beta = 1-f$, so $1-\beta^2 = 1-(1-f)^2 = f(2-f)$.
      Therefore
      \[
        \mathcal{P}_{\tau,x}^{\mathrm{Q5b}}\Sigma^{\mathrm{tot}}(x,n)
        = R(x)\sqrt{f(x,n)(2-f(x,n))}.
      \]
      Equation~\eqref{eq:quadratic-local} follows: $R^2\beta^2 + R^2(1-\beta^2) = R^2$.

      \textbf{Uniqueness.}
      Let $\psi_x(f)$ satisfy (L$_x$), (BI$_x$), (W$_x$).
      From (W$_x$): $B(x,n) = R(x)(1-f)$.
      Substituting into (BI$_x$): $R(x)^2(1-f)^2 + \psi_x(f)^2 = R(x)^2$, giving
      $\psi_x(f) = R(x)\sqrt{f(2-f)}$.
      The positive branch is uniquely forced.

      \textbf{Factorisation.}
      The formula $R(x)\sqrt{f(2-f)}$ separates into a local factor $R(x)$ depending only
      on position and a universal factor $\sqrt{f(2-f)}$ which is identical to the
      homogeneous case of~\cite{Beau2026tp}.
      This factorisation~\eqref{eq:factorisation} is not postulated; it follows from the
      normalised-capacity definition of $\beta$.
    \end{proof}

    \begin{corollary}[Homogeneous case]\label{cor:homogeneous}
      Setting $R(x) \equiv 1$ (vacuum, $\epsilon \equiv 0$) in Theorem~\ref{thm:main}
      recovers the flat-space result of~\cite{Beau2026tp}:
      \[
        \mathcal{P}_\tau^{\mathrm{Q5b}}\Sigma_n^{\mathrm{tot}} = \sqrt{1-B_n^2},
        \qquad
        B_n^2 + \bigl(\mathcal{P}_\tau^{\mathrm{Q5b}}\Sigma_n^{\mathrm{tot}}\bigr)^2 = 1.
      \]
    \end{corollary}

    \begin{proof}
      Immediate from Theorem~\ref{thm:main} with $R\equiv 1$:
      $1\cdot\sqrt{f(2-f)} = \sqrt{(1-B_n)(1+B_n)} = \sqrt{1-B_n^2}$.
    \end{proof}

  \section{Gravitational Time Dilation}\label{sec:timedil}

  \begin{corollary}[Fibrewise Lorentz factor and gravitational time dilation]\label{cor:timedil}
    Under the hypotheses of Theorem~\ref{thm:main}, the fibrewise temporally projected
    saturation ratio is
    \begin{equation}\label{eq:eta-local}
      \eta_x^\tau := \frac{B(x,n)}{\mathcal{P}_{\tau,x}^{\mathrm{Q5b}}\Sigma^{\mathrm{tot}}(x,n)}
      = \frac{R(x)\beta}{ R(x)\sqrt{1-\beta^2}} = \beta\gamma,
    \end{equation}
    independent of $R(x)$, where $\gamma = (1-\beta^2)^{-1/2}$.
    Consequently $\Phi(\eta_x^\tau) = 1/\gamma$ as in the homogeneous case.

    The temporal residual per unit cascade depth at position $x$ is
    \begin{equation}\label{eq:timedilation}
      \frac{d\tau}{dn}\bigg|_x
      = |F^\tau(x,n)| = R(x)\,\frac{1}{\gamma}
      = \frac{R(x)}{\gamma}.
    \end{equation}
    Gravitational time dilation is the factor $R(x) \leq 1$ multiplying the special-relativistic
    dilation $1/\gamma$:
    \begin{equation}\label{eq:ratio}
      \frac{(d\tau/dn)|_x}{(d\tau/dn)|_\infty}
      = \frac{R(x)/\gamma}{1/\gamma} = R(x).
    \end{equation}
    Near a massive body $R(x) < 1$, and clocks at position $x$ run slower than clocks at
    infinity by the factor $R(x)$.
  \end{corollary}

  \begin{proof}
    The ratio $\eta_x^\tau$ in~\eqref{eq:eta-local} cancels $R(x)$ in numerator and
    denominator.
    Equation~\eqref{eq:timedilation} is $|F^\tau| = R(x)\sqrt{1-\beta^2} = R(x)/\gamma$.
    Equation~\eqref{eq:ratio} follows by comparison with the flat-space value $R(\infty) = 1$.
  \end{proof}

  \begin{remark}[Light cone at $x$]\label{rem:lightcone}
    $\beta \to 1$ (maximal spatial mobility) forces $F^\tau \to 0$ for any $R(x) > 0$:
    the light cone condition is local and independent of the gravitational field.
    At $x$ with $R(x) = 0$ (if such a point existed), all temporal capacity is exhausted
    by the local configuration; this would correspond to a horizon-like saturation locus
    or to a breakdown of the effective metric description.
  \end{remark}

  \section{Weak-Field Limit}\label{sec:weakfield}

  \subsection{Connecting $R(x)$ to $g_{\tau\tau}(x)$}

    From Definition~\ref{def:lapse}: $R(x)^2 = -2\,g_{\tau\tau}(x)$.
    The effective co-metric $g^{\mu\nu}(x)$ is determined by the infrared Einstein response
    of~\cite{Beau2026h}: the renormalized spectral entropy variation gives
    $\delta S^\mathrm{ren}_\Pi/\delta g^{\mu\nu} \propto G_{\mu\nu}$, leading to the
    effective Einstein equations.

    In the static weak-field (Newtonian) regime, the Einstein equations give
    \begin{equation}\label{eq:gtt}
      g_{\tau\tau}(x) \;\approx\; -\tfrac{1}{2}\Bigl(1+\frac{2\Phi_N(x)}{c^2}\Bigr),
    \end{equation}
    where $\Phi_N(x)$ is the Newtonian potential satisfying $\nabla^2\Phi_N = 4\pi G_N \rho$.
    The $1/r$ profile of $\Phi_N$ arises as the Green function of the Laplace-type
    operator $A_g$ of~\cite{Beau2026h}: this is the same mechanism as the Newtonian
    interaction emerging from the spectral resolvent (Section~2.3 of~\cite{Beau2026h}).

  \subsection{Weak-field result}

    \begin{corollary}[Weak-field gravitational time dilation]\label{cor:weakfield}
      Under the hypotheses of Theorem~\ref{thm:main} and in the weak-field regime
      $|\Phi_N(x)| \ll c^2$, the local capacity radius satisfies
      \begin{equation}\label{eq:Rweak}
        R(x) = \sqrt{-2\,g_{\tau\tau}(x)} \;\approx\; 1+\frac{\Phi_N(x)}{c^2},
      \end{equation}
      and the gravitational time dilation formula~\eqref{eq:timedilation} becomes
      \begin{equation}\label{eq:timedilation-weak}
        \frac{d\tau}{dn}\bigg|_x
        \;\approx\; \Bigl(1+\frac{\Phi_N(x)}{c^2}\Bigr)\,\frac{1}{\gamma},
      \end{equation}
      in agreement with the standard result of general relativity at leading weak-field order.
    \end{corollary}

    \begin{proof}
      In the Q5b/Q11 convention ($g^{\mu\nu} = 2\eta^{\mu\nu}$), the flat-space inverse metric
      component is $g_{\tau\tau}^\infty = -1/2$.
      The weak-field Einstein equations of~\cite{Beau2026h} give, for a static Newtonian potential,
      $g_{\tau\tau}(x) \approx -\frac{1}{2}\bigl(1 + 2\Phi_N(x)/c^2\bigr)$.
      Substituting into~\eqref{eq:lapse}:
      \[
        R(x)^2 = -2\,g_{\tau\tau}(x)
        = -2 \cdot \Bigl(-\tfrac{1}{2}\Bigr)\bigl(1+2\Phi_N/c^2\bigr)
        = 1 + \frac{2\Phi_N(x)}{c^2}.
      \]
      At first order in $\Phi_N/c^2$: $R(x) \approx 1 + \Phi_N(x)/c^2$.
      Equation~\eqref{eq:timedilation-weak} follows from~\eqref{eq:timedilation}.
      The 1/r profile of $\Phi_N(x)$ is the Green function of $A_g$ in three spatial dimensions,
      as established in~\cite{Beau2026h} Section~2.3.
    \end{proof}

    \begin{remark}[Status of the weak-field result]\label{rem:status-weak}
      Corollary~\ref{cor:weakfield} is conditional on two imports:
      (i) the effective Einstein equations from~\cite{Beau2026h} (in the infrared,
      weak-curvature, regime $R\ell^2_{\mathrm{sp}} \ll 1$);
      (ii) the identification of $\Phi_N(x)$ as the Newton coupling via the Green
      function of $A_g$ (Section~2.3 of~\cite{Beau2026h}).
      No new hypothesis is introduced here.
      The result is a consequence of combining the capacity-lapse language with the
      gravitational output of~\cite{Beau2026h}.
    \end{remark}

  \section{Logical Status and Open Problems}\label{sec:status}

  \subsection{Status table}

    \begin{center}
      \renewcommand{\arraystretch}{1.2}
      \begin{tabular}{lll}
        \hline
        Label & Content & Status \\
        \hline
        \cite{Beau2026tp} Thm.~4.5 & Homogeneous map $\sqrt{1-B_n^2}$, uniqueness &
        Proved (under [U]) \\
        Q11~\cite{Beau2026q11} & $A_\tau = 2$, flat metric $g^{\mu\nu} = 2\eta^{\mu\nu}$ &
        Proved (under [U]) \\
        Definition~\ref{def:radius} & $R(x)^2 = 1-\epsilon(x)$ &
        Structural (this paper) \\
        Definition~\ref{def:lapse} & $R(x)^2 = -2g_{\tau\tau}(x)$ &
        Structural (this paper) \\
        Theorem~\ref{thm:main} & Fibrewise map, uniqueness &
        Proved (under [U] + fibrewise BI lift) \\
        Corollary~\ref{cor:timedil} & $d\tau/dn|_x = R(x)/\gamma$ &
        Proved (under [U] + fibrewise BI lift) \\
        Corollary~\ref{cor:weakfield} & $R(x) \approx 1+\Phi_N/c^2$ &
        Conditional (imports~\cite{Beau2026h}) \\
        \hline
      \end{tabular}
    \end{center}

  \subsection{What is not proved here}

    The present paper does not derive the Einstein field equations.
    That derivation is in~\cite{Beau2026h}.
    The present paper does not derive $R(x)$ from first principles as a function of the
    spectral occupancy $\epsilon(x)$.
    It gives two definitions (Definitions~\ref{def:radius} and~\ref{def:lapse}) that are
    equivalent in the Einstein regime, but the bridge $\epsilon(x) \leftrightarrow g_{\tau\tau}(x)$
    is not closed analytically within the capacity-flow language alone.
    Closing this bridge is LC-O2's remaining open problem (LC-O2-O1 below).

  \subsection{Open problems}

    \paragraph{LC-O2-O1 (Spectral closure of $R(x) \leftrightarrow g_{\tau\tau}(x)$).}
      Definitions~\ref{def:radius} and~\ref{def:lapse} give two expressions for $R(x)$:
      one from spectral occupancy $\epsilon(x)$ and one from the effective metric.
      A direct spectral derivation showing
      \[
        \epsilon(x) = 1-R(x)^2 = 1+2g_{\tau\tau}(x)
      \]
      from the admissibility data alone would close this identification without importing
      the Einstein equations.
      This is the capacity-flow analogue of deriving gravitational time dilation from the
      spectral operator structure directly.

    \paragraph{LC-O2-O2 (Length contraction in the fibrewise regime).}
      Corollary~\ref{cor:timedil} gives fibrewise time dilation and the light cone.
      Length contraction in the non-homogeneous setting requires a comparison of
      spatial capacity across different fibres and the reconstruction of effective Lorentz
      transformations between locally inertial observers in the curved geometry.

    \paragraph{LC-O2-O3 (Strong-field regime: $R(x) \to 0$).}
      Remark~\ref{rem:lightcone} identifies $R(x) = 0$ with a horizon-like saturation locus
      or a breakdown of the effective metric description.
      The capacity-flow language provides a novel criterion for horizon formation:
      the local capacity budget is fully saturated by the configuration.
      Whether this criterion reproduces the Schwarzschild horizon condition is open.

  \section*{Acknowledgements}

  The author acknowledges the use of large language models as a supportive tool for refining
  language, structure, and internal consistency during the development of this manuscript.

  \newpage

  \bibliographystyle{plain}
  \bibliography{cosmochrony-bibliography}

\end{document}
